\section{Introduction}

Photodiode TIAs convert photocurrent into voltage and often sit at the front of optical receiver or detector readout chains. Their apparent simplicity hides a coupled design space. Feedback resistance sets nominal low-frequency transimpedance, while photodiode capacitance, input capacitance, feedback capacitance, and finite op-amp bandwidth shape the closed-loop response. A design point that preserves bandwidth can exhibit substantial peaking under some assumptions, while a more conservative compensation choice can reduce peaking at the cost of bandwidth. Published TIA work treats photodiode capacitance, bandwidth extension, regulated-cascode input stages, feedback structures, and noise as recurring design concerns \citep{ANALUI_2004_BW_ENHANCEMENT,LU_2007_BROADBAND_TIA,PARK_YOO_2004_RGC_TIA,NOH_2020_CF_TIA_DC_LOOP}.

This project addresses the early screening stage of photodiode TIA design. Instead of presenting a new circuit topology, it builds a traceable workflow for exploring selected design assumptions before detailed vendor simulation, layout-aware analysis, or hardware work. The workflow produces MATLAB-generated response curves, sweep summaries, design-region maps, vendor macromodel comparison figures, and first-pass behavioural noise estimates. These artifacts support cautious manuscript claims because each figure and table is tied to committed scripts, CSV files, and limitation notes.

Finite op-amp bandwidth is central to the workflow. A single-pole behavioural approximation makes the relationship between $R_f$, $C_f$, $C_{pd}$, $A_0$, and $f_t$ explicit enough for fast parameter sweeps. The same simplification also defines the model boundary: the behavioural model does not capture all vendor pole-zero behaviour, output limits, nonlinearities, package effects, or layout parasitics. The baseline behavioural magnitude and phase responses provide the first visual anchor for this finite-bandwidth model (Figure~\ref{fig:baseline-response}). A compact model-parameter table is provided in the methodology section to make the assumed and swept quantities explicit (Table~\ref{tab:model-assumptions}).

The resulting contribution is a reproducible and claim-bounded pre-design workflow. It maps selected operating assumptions, visualizes project-defined Safe/Marginal/Risky regions, compares selected trends with vendor macromodel data, and adds first-pass noise context. The workflow helps organize early-stage design evidence; it does not establish a universal TIA design rule.
